\section{Discussion}
\label{sec:discussion}

\subsection{Interpretation of the Verification Evidence}
\label{subsec:discussion-verification-evidence}

The verification evidence supports a bounded implementation claim: the reported
solver configuration, benchmark records, and output summaries reproduce the
finite-volume and comparison records used in this report. The implementation
checks exercise descriptor health, sampled self-convergence, native/SciML
cross-integrator discrepancies, stationary-Stokes projection rows, boundary
handling, rheology/profile sensitivity, and resolved-velocity output
diagnostics. These checks make the numerical record traceable rather than
anecdotal.

The same evidence should not be read as a stronger theorem or validation claim.
The benchmark rows do not prove convergence of every model component, do not
establish pressure accuracy, and do not show agreement with clinical or
experimental measurements. For the current comparison cases, the generated
summaries persist characteristic-speed diagnostics; CFL, mass, and positivity
diagnostics are exposed by the solver and verification records rather than as
tracked C23/C40 comparison time histories. The remaining limitations are that
the zero-forcing rest-state test reveals a material non-well-balanced start-up
artifact, and production grid, timestep, start-up, and stationarity sensitivity
remain incomplete for the C23/C40 comparison runs. The present result is
therefore best described as a checked implementation record with documented
limitations, not as complete numerical verification of every possible output.

\subsection{Interpretation of the Cross-Model Discrepancies}
\label{subsec:discussion-cross-model-discrepancies}

The largest resolved-velocity disagreement in this pair occurs in C40, the
stronger stenosis case. The section-mean comparison in
Table~\ref{tab:t1-3d-comparison-summary} records the descriptive widening:
the velocity $L^1$ error increases from 0.505 cm/s in C23 to 0.966 cm/s in C40,
while the velocity $L^3$ error increases from 0.853 cm/s to 2.82 cm/s. The
largest section discrepancy increases from 2.35 cm/s to 9.92 cm/s.
Table~\ref{tab:t1-largest-section-discrepancies} localizes the largest section
discrepancies just upstream of the displayed throat band, in the C40
acceleration/shoulder region. Because the numerical and data-matching gates
remain unresolved, this two-case pattern should not be read as a causal
stenosis-severity result.

The radial-profile comparison gives a complementary reading. The throat and
downstream C40 slices show the largest profile discrepancies, especially in
Table~\ref{tab:t1-radial-profile-summary}. That is consistent with the
expectation that a one-dimensional parabolic reconstruction will be most
strained where the resolved field contains strong transverse structure. It is
not, by itself, a diagnosis of secondary flow, skewness, separation, or any one
omitted 3D mechanism. The radial plots show where the 1D profile closure fails
to reproduce the resolved radial bins; they do not identify a unique physical
cause.

\subsection{Effects of Unmatched Wall Models, Boundary Conditions, and Histories}
\label{subsec:discussion-unmatched-modeling-effects}

The present comparison changes more than spatial dimension. The 1D computation
uses a compliant reduced Canic--Koiter wall law and a fixed-area characteristic
boundary approximation, while the local resolved files expose velocity,
pressure, and displacement as single final-time snapshots rather than a matched
wall-history record. These differences matter because velocity
discrepancy is not solely a profile-closure discrepancy. Area mismatch can change
$\bar u=Q/A$ even when physical flow is close; flow mismatch can arise from
different inlet, outlet, wall, or start-up histories; and the final-time
difference between $1.0$ s and $0.9995$ s does not establish matched phases or
matched histories.

The cut-area audit in Table~\ref{tab:t1-area-audit} is therefore an operator
check, not an attribution result. It shows typical sub-percent agreement
between the tetrahedral cuts and the static reference areas, with localized
multi-percent outliers. Table~\ref{tab:t1-flow-area-decomposition} adds a
section-wise split between physical-flow and area-denominator contributions for
the reported samples; that algebraic split is useful bookkeeping, but it does
not by itself establish whether the flow mismatch is physical, numerical, or
data-history driven. The comparison should therefore be read as a coupled
velocity, flow, geometry, and modeling diagnostic rather than as a pure test of
the cross-sectional profile closure.

\subsection{Methodological Limitations}
\label{subsec:discussion-methodological-limitations}

The most important limitation is evidential scope. The report compares a
selected 1D realization with two resolved computational velocity datasets under
a declared plane--tetrahedron output operator. It does not compare pressures,
pressure drops, clinical FFR, CT-FFR, wall shear stress, structural deformation,
or experimental measurements. The pressure-ratio notation in the report
therefore remains an output convention for later studies, not a completed
pressure-ratio result.

A second limitation is that the comparison is not a fully matched 1D--3D
validation study. Matching would require, at minimum, a declared common
geometry, wall model, rheology, material parameters, inlet and outlet histories,
initialization, time alignment, numerical-resolution study, and output
operator. Some of those ingredients are available here, especially the analytic
reference geometry and the velocity output operator. Others are absent or only
partially persisted, so the current comparison can localize discrepancies but
cannot uniquely attribute them. The section-flow audit is another unresolved
gate: the available 3D cuts vary by about 7.7\% of the C23 mean section flow
and 20.2\% of the C40 mean section flow, and the comparison files do not by
themselves establish whether this is fixed-wall numerical/operator error,
single-snapshot moving-wall geometry, or a moving-wall effect requiring
area-rate and wall-history data.

\subsection{Implications and Future Work}
\label{subsec:discussion-implications-future-work}

The immediate implication is practical: the next useful step is a more
controlled comparison record. That record should persist matched inlet and
outlet histories, wall and material parameters, positivity/CFL/mass diagnostics,
subcritical characteristic-speed signs, grid and operator sensitivity for the
actual C23/C40 outputs, high-area-error section sensitivity, and an explicit
classical-versus-extended ablation. It should also preserve the existing
physical-flow and area-denominator decomposition as a standard output rather
than as an after-the-fact table.

A pressure-focused extension would require a separate design. It would need
proximal and distal pressure observation maps, reference-pressure and averaging
conventions, pressure data from the comparison model, and a clear statement of
whether the target is a reduced pressure diagnostic, a pressure-drop comparison,
or a clinical pressure-ratio comparison. Until such a study is added, the
defensible contribution remains the audited extended-1D model specification,
the reproducible numerical workflow, and the declared operator for comparing 1D
and resolved-3D velocity outputs.

\section{Conclusions and Limitations}
\label{sec:conclusions-limitations}

This report has separated three claims that are often blurred in reduced-order
hemodynamics work: model specification, numerical verification, and cross-model
comparison. The selected model is now stated in both physical variables
$(A_{\mathrm{phys}},Q_{\mathrm{phys}})$ and solver variables
$(a,q)=(A_{\mathrm{phys}}/\pi,Q_{\mathrm{phys}}/\pi)$. It is an
$R_{\max}$-normalized Canic-derived, radius-squared finite-volume model with
$R_0^*=R_{\max}$ in the continuous source-balanced elastic wall law, a
parabolic-profile baseline closure, optional Canic variable-radius terms, and a
fixed-area characteristic boundary approximation.

The verification record documents an internally reproducible numerical record
for the selected solver and recorded artifacts. Descriptor checks,
self-convergence rows, cross-integrator discrepancies, stationary-Stokes
projection rows, boundary diagnostics, rheology/profile sensitivity,
and resolved-velocity output diagnostics support the
reported finite-volume/time-integration data contracts. They remain summary
diagnostics rather than full time histories, and the corrected rest-state test
shows that the selected discretization is not well balanced for the stenotic
geometry-rest state. At the finest checked grid, the zero-input artificial flow
at $t=1.0$ s is still of the same order as the comparison-flow scale, so the
later C23/C40 velocity comparison cannot be interpreted as a production-accuracy
study.

The 1D--3D velocity comparison shows localized discrepancy under the declared
cross-section quadrature operator. For C23 and C40 at $t=1.0$ s, the
section-wise area-mean velocity discrepancies have $L^1$ errors of 0.505 and
0.966 cm/s, $L^2$ errors of 0.664 and 1.87 cm/s, and $L^3$ errors of 0.853 and
2.82 cm/s. The largest listed section discrepancies are 2.35 and 9.92 cm/s.
The larger discrepancies in this two-case pair occur in the upstream C40
shoulder and in the reconstructed radial profiles; they are single-realization
descriptive comparison outputs, not a validated stenosis-dependence or
model-accuracy claim.

The bounded contribution is therefore precise rather than clinical: a selected
extended 1D stenosis model, an auditable finite-volume/time-integration
implementation, and a cross-sectional velocity comparison operator that identifies
where the current 1D and resolved-3D computational outputs differ. The report
does not establish pressure-drop accuracy, pressure-ratio or FFR behavior,
clinical validity, full transient 3D parity, structural deformation, or a fully
matched 1D--3D validation. Those remain future studies with different evidence
requirements.
